% ======================================================================
% Chapter 1: Introduction
% Source: R. Gupta, "Introduction to Lattice QCD" (arXiv:hep-lat/9807028),
% Section 1, pp. 5-7. The closing paragraphs continue onto page 7 (their
% text is extracted at the top of sec02.txt).
% ======================================================================

\chapter{Introduction}

The goal of the lectures on lattice QCD (LQCD) is to provide an overview
of both the technical issues and the progress made so far in obtaining
phenomenologically useful numbers. The lectures consist of three parts.
My charter is to provide an introduction to LQCD and outline the scope
of LQCD calculations. In the second set of lectures, Guido Martinelli will
discuss the progress we have made so far in obtaining results, and their
impact on Standard Model phenomenology. Finally, Martin L\"{u}scher will
discuss the topical subjects of chiral symmetry, improved formulation of
lattice QCD, and the impact these improvements will have on the quality
of results expected from the next generation of simulations.

QCD is the regnant theory of strong interactions. It is formulated in
terms of quarks and gluons which we believe are the basic degrees of
freedom that make up hadronic matter. It has been very successful in
predicting phenomena involving large momentum transfer. In this regime the
coupling constant is small and perturbation theory becomes a reliable tool.
On the other hand, at the scale of the hadronic world, $\mu \lesssim
1~\mathrm{GeV}$, the coupling constant is of order unity and perturbative
methods fail. In this domain lattice QCD provides a non-perturbative tool
for calculating the hadronic spectrum and the matrix elements of any
operator within these hadronic states from first principles. LQCD can also
be used to address issues like the mechanism for confinement and chiral
symmetry breaking, the role of topology, and the equilibrium properties of
QCD at finite temperature. Unfortunately, these latter topics were not
covered at this school, so I will give appropriate references to compensate
for their omission.

Lattice QCD is QCD formulated on a discrete Euclidean space time grid.
Since no new parameters or field variables are introduced in this
discretization, LQCD retains the fundamental character of QCD. Lattice QCD
can serve two purposes. First, the discrete space-time lattice acts as a
non-perturbative regularization scheme. At finite values of the lattice
spacing $a$, which provides an ultraviolet cutoff at $\pi/a$, there are no
infinities. Furthermore, renormalized physical quantities have a finite
well behaved limit as $a \to 0$. Thus, in principle, one could do all the
standard perturbative calculations using lattice regularization, however,
these calculations are far more complicated and have no advantage over
those done in a continuum scheme. Second, the pre-eminent use of
transcribing QCD on to a space-time lattice is that LQCD can be simulated
on the computer using methods analogous to those used for Statistical
Mechanics systems. (A brief review of the connection between Euclidean
field theory and Statistical Mechanics is given in Section~5.) These
simulations allow us to calculate correlation functions of hadronic
operators and matrix elements of any operator between hadronic states in
terms of the fundamental quark and gluon degrees of freedom.

The only tunable input parameters in these simulations are the strong
coupling constant and the bare masses of the quarks. Our belief is that
these parameters are prescribed by some yet more fundamental underlying
theory, however, within the context of the standard model they have to be
fixed in terms of an equal number of experimental quantities. This is what
is done in LQCD. Thereafter all predictions of LQCD have to match
experimental data if QCD is the correct theory of strong interactions.

A very useful feature of LQCD is that one can dial the input parameters.
Therefore, in addition to testing QCD we can make detailed predictions of
the dependence of quantities on $\alpha_s$ and the quark masses. These
predictions can then be used to constrain effective theories like chiral
perturbation theory, heavy quark effective theory, and various
phenomenological models.

My first lecture will be devoted to an overview of the scope of LQCD and
to showing that simulations of LQCD are a step by step implementation of
field theory. The second lecture will be devoted to explaining the details
of how to transcribe the quark and gluon degrees of freedom on to the
lattice, and to construct an action that, in the limit of zero lattice
spacing, gives continuum QCD. I will also spend some time on issues of
gauge invariance, chiral symmetry, fermion doubling problem, designing
improved actions, the measure of integration, and gauge fixing.

Numerical simulations of LQCD are based on a Monte Carlo integration of
the Euclidean path integral, consequently, the measurements have
statistical errors in addition to the systematic errors due to lattice
discretization. In order to judge the quality of lattice calculations it is
important to understand the origin of these errors, what is being done to
quantify them, and finally what will it take to achieve results with a
given precision. These issues will be covered in the third lecture along
with an elementary discussion of Monte Carlo methods.

The fourth lecture will be devoted to the most basic applications of
LQCD---the calculation of the hadron spectrum and the extraction of quark
masses and $\alpha_s$. Progress in LQCD has required a combination of
improvements in formulation, numerical techniques, and in computer
technology. My overall message is that current LQCD results, obtained from
simulations at $a \approx 0.05$--$0.1$~fermi and with quark masses roughly
equal to $m_s$, have already made an impact on Standard Model
phenomenology. I hope that the three sets of lectures succeed in
communicating the excitement of the practitioners.

In a short course like this I can only give a very brief description of
LQCD. In preparing these lectures I have found the following books
[1--4], and reviews [5--8] very useful. In addition, the proceedings of
the yearly conferences on lattice field theory [9] are excellent sources
of information on current activity and on the status of results. I hope
that the lectures at this school and these references will allow any
interested reader to master the subject.
